\item En una muestra de un metal sólido, cada átomo tiene libertad de vibrar en torno a alguna posición de equilibrio. La energía del átomo consiste en energía cinética para movimiento en las direcciones $x, y$ y $z$, más energía potencial elástica asociada con las fuerzas de la ley de Hooke ejercidas por los átomos vecinos sobre él en las direcciones $x, y$ y $z$. De acuerdo con el teorema de equipartición de la energía, suponga que la energía promedio de cada átomo es $\frac{1}{2}k_BT$ para cada grado de libertad. (a) Pruebe que el calor específico molar del sólido es $3R$. La ley de Dulong-Petit establece que este resultado describe sólidos puros a temperaturas suficientemente altas. (Puede ignorar la diferencia entre el calor específico a presión constante y el calor específico a volumen constante.)
